%% parts/p12/first-order-automation.tex --- Part 12.3
\chapter{First-Order Automation}
\label{sec:p12-first-order-automation}

\section{Canonical unification and resolution interfaces}

The exact substitution, instance, unifier, most-general-unifier, and elementary
unification definitions, together with the source's unification result and
proof sketch, occur in Section~\ref{sec:fol-cardinality-normalforms-04}.
Those particular definitions and that result are unlabeled in the exact
source.  P12 therefore does not invent labels for the unlabeled source units
or rename the proof sketch as a complete Martelli--Montanari proof.

\begin{hbexample}{def:p12-occurs-check}{The finite-term occurs-check boundary}
Unifying \(P(f(x,x))\) with \(P(f(g(y),y))\) requires both \(x=g(y)\) and
\(x=y\).  Eliminating \(x\) leaves \(y=g(y)\), so the variable would have to
occur properly inside its own finite image.  The occurs check therefore
rejects the equation, and no finite first-order term substitution unifies the
atoms.  Changing the term domain to permit such nonfinite self-referential
solutions is outside the canonical P2 theorem.
\end{hbexample}

\begin{hbmethod}{def:p12-resolution-search}{Resolution search with an auditable derivation record}
\index{resolution search}
This method name is needed to distinguish clause-selection and redundancy
control from the inference rules themselves.  Begin with the exact finite
clause set and record every signature-changing clausification step.  Before a
binary inference, rename the parent occurrences apart as specified in
Section~\ref{sec:fol-cardinality-normalforms-04}; apply the declared refinement's
literal restriction; call the P2 most-general-unifier
procedure; and record the resulting instance of
\hyperref[sheet:resolution-fol]{\texttt{sheet:resolution-fol}}.  Delete or
simplify clauses only under a declared redundancy operation, and retain a DAG
whose leaves are input clauses and whose internal nodes name their sheet
instances.  Thus the method stands as a search-and-record discipline over a
registered calculus, not as an additional inference rule.  Its output is used
by the soundness/completeness audit below.
\end{hbmethod}

Resolution and Factoring soundness and its complete proof occur inline in
Section~\ref{sec:fol-cardinality-normalforms-04}; the exact source theorem is
unlabeled.  The separately identified editorial lookup
\hyperref[sheet:resolution-fol]{\texttt{sheet:resolution-fol}} provides
navigation and comparison but does not replace the inherited Rule displays.
The equality-free, finite-clause refutational-completeness result has its exact
inherited proof sketch at Theorem~\ref{thm:fol-resolution}, including
its fresh-ground-term and factoring boundary.  Figure~\ref{fig:p02-fol-resolution-dag}
is the first-order instance of the resolution-refutation diagram; the propositional master is
Figure~\ref{fig:p02-resolution-dag}.

\begin{hbexample}{def:p12-clausification-audit}{Equivalence before Skolem equisatisfiability}
For the negation of
\(\forall x(P(x)\to\exists y\,R(x,y))\), ordinary Boolean and quantifier
rewrites give the formula-equivalence chain
\[
 \neg\forall x(P(x)\to\exists y\,R(x,y))
 \ \Leftrightarrow\
 \exists x\bigl(P(x)\land\forall y\,\neg R(x,y)\bigr).
\]
Introducing a fresh constant \(c\) at the existential step yields the clauses
\(\{P(c)\}\) and \(\{\neg R(c,y)\}\).  Their universally closed
interpretations are equisatisfiable with the negated sentence in the expanded
signature; they are not asserted to be the same-signature formula equivalent.
This example routes the full pipeline to the P2 normal-form treatment rather
than duplicating it.
\end{hbexample}

\section{Resolution refinements}

\begin{hbdefinition}{def:p12-resolution-refinement-roles}{Refinement roles}
A \emph{resolution refinement} is a precisely declared restriction, inference
family, search strategy, or redundancy operation applied around a fixed base
calculus.  The umbrella name is necessary for comparison but does not erase
those four different standings.  For the ordered-resolution surface compared
here, fix a term ordering that is well founded, stable under substitution,
monotone under term contexts, and total on ground terms, together with its
declared extension to literals.  A \emph{literal selection function} maps each
clause to a subset of its negative literals.  A clause literal is
\emph{eligible} when it is selected, or, if that clause has no selected
literal, when it is maximal in the declared literal ordering; ordered
resolution admits only inferences whose resolved literals are eligible after
the unifier is applied.  Thus the ordering restriction and literal selection
are explicit and neither means that the underlying Resolution rule changed.

In Robinson's nucleus/electron
notation for hyperresolution, a nucleus clause \(N\) has designated literals
that are paired with designated complementary literals in electron clauses
\(E_1,\ldots,E_q\); one simultaneous most general unifier makes all pairs
complementary, and the hyperresolvent contains the unpaired literals of all
premises under that unifier.  A concrete positive or negative variant must
also fix the required polarities and factoring convention.

In Slagle's semantic-resolution notation, fix a ground interpretation \(M\).
A \emph{semantic clash} is a finite family
\(\{E_1,\ldots,E_q,N\}\) in which every electron \(E_i\) is false under
\(M\), \(R_1=N\), each \(R_{i+1}\) is an ordinary resolvent of \(R_i\) and
\(E_i\), and the final \(R_{q+1}\) is false under \(M\).  The clause \(N\)
is the nucleus and \(R_{q+1}\) is the \(M\)-resolvent.  This clash condition
is the exact fixed-interpretation eligibility test used here: it restricts
which multi-premise resolution chains are admitted, while every conclusion
remains a syntactic resolution consequence.  It differs from hyperresolution's
polarity restriction and from the support partition below
\cite{Robinson1965,Slagle1967}.

For \emph{set of support}, partition the input clauses into a usable set
\(U\) and a support set \(S\).  The support closure is the least clause set
containing \(S\) and containing the conclusion of every Resolution or
Factoring instance on
\hyperref[sheet:resolution-fol]{\texttt{sheet:resolution-fol}} that has at
least one premise already in that closure and has every other premise in
\(U\) or that closure.  The strategy permits exactly those instances.  With
no deletion or contraction, such a search is \emph{fair} when every permitted
sheet instance is eventually performed after all of its premises have been
generated.  This is a parent-selection strategy, not a new
inference family.

A clause \(C\) \emph{subsumes} \(D\) exactly
when some term substitution \(\rho\) satisfies \(C\rho\subseteq D\), licensing
deletion of \(D\) as redundant.  None is synonymous with unrestricted
resolution, and each later completeness claim needs its own hypotheses.
\index{ordered resolution}\index{hyperresolution}\index{hyperresolvent}
\index{semantic resolution}
\index{literal selection function}\index{eligible literal}
\index{semantic clash}\index{nucleus}\index{electron}
\index{M-resolvent}
\index{set of support}\index{support closure}\index{fair search}
\index{subsumption}
\end{hbdefinition}

\begin{longtable}{@{}p{.16\textwidth}p{.25\textwidth}p{.29\textwidth}p{.20\textwidth}@{}}
\caption{Resolution-refinement taxonomy.}
\label{tab:p12-resolution-refinements}\\
\toprule
Refinement & Standing & Completeness/soundness boundary & Evidence \\ \midrule
\endfirsthead
\caption[]{Resolution-refinement taxonomy (continued)}\\
\toprule
Refinement & Standing & Completeness/soundness boundary & Evidence \\ \midrule
\endhead
\midrule\multicolumn{4}{r}{\emph{continued on next page}}\\
\endfoot
\endlastfoot
Ordered resolution & The ordering-and-literal-selection eligibility
restriction defined above & Completeness requires the exact admissibility,
redundancy, and fair-saturation hypotheses of the cited presentation;
unrestricted completeness does not transfer automatically &
\cite{BachmairGanzinger2001} \\ \addlinespace
Hyperresolution & The nucleus/electron multi-premise inference family defined
above & Its rule surface and completeness theorem must fix polarity and
factoring; it is not identified with an informal batch of binary steps &
\cite{Robinson1965} \\ \addlinespace
Semantic resolution & Slagle's \(M\)-clash restriction with false electrons
and a false final \(M\)-resolvent & The fixed interpretation determines clash
eligibility; derived clauses remain syntactic consequences, and the cited
completeness result belongs to this exact restriction &
\cite{Slagle1967} \\ \addlinespace
Set of support & The support-closure parent restriction defined above &
Completeness requires a satisfiable usable partition and fair search, as stated
in Theorem~\ref{thm:p12-set-of-support} &
\cite{WosRobinsonCarson1965} \\ \addlinespace
Subsumption & Contraction/redundancy deletion & Sound deletion requires the
instance-containment condition; it creates no new consequence and its use must
preserve the declared saturation fairness & \cite{BachmairGanzinger2001} \\
\bottomrule
\end{longtable}

\begin{hbtheorem}{thm:p12-set-of-support}{cited=WosRobinsonCarson1965}{Set-of-support completeness boundary}
Partition an equality-free first-order clause set into a usable set \(U\) and
a support set \(S\), and use the support closure and fair search of
Definition~\ref{def:p12-resolution-refinement-roles}, without clause deletion
or contraction.  If \(U\) is satisfiable while \(U\cup S\) is unsatisfiable,
every such fair search derives the empty clause.  Dropping the
satisfiability of \(U\), factoring, or fairness is not licensed by this
statement.
\end{hbtheorem}

\section{Equality and the deferred exact calculi}

\begin{proofboundary}{def:p12-resolution-equality-boundary}{Equality is not an uninterpreted predicate}
Ordinary first-order Resolution and Factoring are sound when equality atoms are
treated like other atoms, but their cited refutational-completeness theorem is
equality-free.  Treating the equality glyph as an arbitrary binary relation
changes identity semantics; adding identity axioms without a named exact
equality calculus changes the search and completeness hypotheses.  Ordered
superposition controls equality replacement through term ordering,
eligibility, and redundancy, none of which is optional implementation detail
\cite{BachmairGanzinger2001}.
\end{proofboundary}

\begin{hbmethod}{def:p12-martelli-montanari}{Martelli--Montanari unification}
The exact transformation system is
\hyperref[sheet:unification-first-order]{\texttt{sheet:unification-first-order}}.
Its retained-equation state has exactly the same unifiers as the initial
multiset, as proved in \cref{thm:p12-ue-sheet-invariant}. Delete and Decompose
preserve syntactic solutions; Orient changes only equation order. Eliminate
retains the binding \(x\doteq t\) and substitutes it in the other equations;
this sheet has no separate substitution accumulator. Clash is unsatisfiable by freeness of the
term algebra, and Occurs fails because no finite term can equal a proper term
containing itself. The method name is necessary to distinguish this finite
syntactic procedure from unification modulo equations or higher-order
unification. It stands as a transformation algorithm, differs through the
occurs check and free-term semantics, and is used in the MGU theorem
\cite{MartelliMontanari1982,BaaderSnyder2001}.
\index{Martelli--Montanari algorithm}\index{occurs check}
\end{hbmethod}

The sheet states no rule priority. Termination under every legal enabled-rule
choice uses the exact isolated-binding and term-node measure in
\cref{def:p12-ue-measure,thm:p12-ue-sheet-termination}. A successful
irreducible state has solved equations \(x_i\doteq t_i\) with distinct
left variables absent from all right sides; simultaneous extraction and its
most-generality argument are \cref{thm:p12-ue-extraction}.

\begin{hbtheorem}{thm:p12-unification-mgu}{cited=MartelliMontanari1982}{Most-general-unifier factorization}
For every finite first-order syntactic unification problem, the displayed
algorithm terminates with failure exactly when no unifier exists. On success it
returns a unifier \(\mu\) such that every unifier \(\theta\) factors as
\(\theta=\delta\circ\mu\) on the variables of the problem, for some
substitution \(\delta\). The factor substitution need not be unique.
The uniqueness-up-to-renaming property concerns most general solutions
themselves under the stated substitution equality, as proved in
\cref{thm:p12-ue-renaming}; it is not a uniqueness claim about every
factor witness. The full literal-sheet and separate accumulator derivations
are \cref{thm:p12-ue-extraction,thm:p12-ue-accumulator}.
\end{hbtheorem}

Associative, commutative, or AC-unification changes equality from syntactic
identity to congruence modulo a theory and may require multiple incomparable
members of a minimal complete set of unifiers, rather than one MGU covering
all solutions.
Higher-order unification permits functional variables and is undecidable in
general \cite{Goldfarb1981}. Neither boundary is implemented by the first-order
sheet.

\begin{hbdefinition}{def:p12-ordered-superposition}{Ordered-superposition saturation}
The calculus is
\hyperref[sheet:ordered-superposition]{\texttt{sheet:ordered-superposition}}.
A clause is redundant in a set \(S\) when it is entailed by strictly smaller
clauses under the fixed ground-clause ordering in the declared redundancy
criterion. A fair saturation derivation eventually performs every persistent
nonredundant eligible inference or makes its conclusion redundant. These names
are needed because unrestricted equality replacement is neither terminating
search nor the cited completeness calculus. They stand as an inference system,
deletion criterion, and scheduling property, differ through the explicit
ordering/selection constraints, and are used in the completeness boundary
\cite{BachmairGanzinger2001,NieuwenhuisRubio2001}.
\index{ordered superposition}\index{saturation}\index{redundancy!superposition}
\end{hbdefinition}

\begin{hbtheorem}{thm:p12-superposition-completeness}{cited=BachmairGanzinger2001}{Refutational completeness of ordered superposition}
For clauses with identity semantics, a complete simplification ordering,
admissible negative-literal selection, the sheet's eligible inferences, a
sound redundancy criterion, and fair saturation, an unsatisfiable input
derives the empty clause. The model-construction proof reads a saturated set
without the empty clause into a rewrite interpretation; deleting any listed
hypothesis changes the theorem.
\end{hbtheorem}

The description-logic obligation is now supplied at its Part-IX home:
Definition~\ref{def:p09-alc} fixes \(\mathcal{ALC}\),
Method~\ref{def:p09-alc-tableau} gives its blocking tableau, and
Theorem~\ref{thm:p09-alc-modal} states the exact multimodal translation. The
classical propositional and first-order tableau sheets remain distinct
calculi.

\section{Practical status of unrestricted first-order search}

Under the fixed effective coding, unrestricted first-order validity is
computably enumerable but undecidable
(Theorem~\ref{thm:fol-validity-undecidable} and its immediately following
enumerability corollary;
\cite{Church1936,Turing1937}).  A sound complete proof search may therefore
run forever on an invalid input; a timeout is not a countermodel.  The statement
does not transfer without proof to every fixed signature or to the decidable
fragments catalogued in Part~2.
